\par\Needspace{6\baselineskip}\medskip
\hypertarget{native-083-article064}{}\pdfbookmark[2]{64.}{article-064}
\begin{center}64.\end{center}
Circa problema art. 62 sequentes adhuc observationes adiicimus.

I. Statuendo in aequatione secunda illic tradita $N=\lambda$, $N=L$, $N=l$, prodit $R'\sin(\lambda-L)=\Delta'\sin(l-\lambda)$; $r'\sin(\lambda-L)=\Delta'\sin(l-L)$; $r'\sin(l-\lambda)=R'\sin(l-L)$; aequatio prima aut secunda commode ad calculi confirmationem applicatur, si methodus I aut II art. 62 adhibita est. Ita habetur in exemplo nostro
\begin{center}
\begin{tabular}{@{}ll@{}}
\makebox[13em][l]{$\log\sin(\lambda-L)$\hfill .\,.\,.} & $9{,}475\,8653_n$\\[.3em]
\makebox[13em][l]{$\log\frac{\Delta'}{r'}$\hfill .\,.\,.\,.\,.\,.\,.\,.\,.} & $9{,}754\,6117$\\[.3em]
\hline
& $9{,}721\,2536_n$\\[.3em]
\makebox[13em][l]{$\log\sin(l-L)$\hfill .\,.\,.} & $9{,}721\,2536_n$
\end{tabular}
\end{center}
\begin{center}
$l-L=-31^0\,45^{\prime}\,26^{\prime\prime}{,}82$
\end{center}

II. Sol duoque in plano eclipticae puncta, quae sunt proiectiones loci corporis coelestis atque loci terrae, triangulum planum formant, cuius latera sunt $\Delta'$, $R'$, $r'$, angulique oppositi vel $\lambda-L$, $l-\lambda$, $180^0-l+L$, vel $L-\lambda$, $\lambda-l$, $180^0-L+l$: ex hoc principio relationes in I traditae sponte sequuntur.

III. Sol, locus verus corporis coelestis in spatio, locusque verus terrae aliud triangulum formabunt, cuius latera erunt $\Delta$, $R$, $r$: angulis itaque his resp. oppositis per $S$, $T$, $180^0-S-T$ denotatis, erit $\frac{\sin S}{\Delta}=\frac{\sin T}{R}=\frac{\sin(S+T)}{r}\cdot$ Planum huius trianguli in sphaera coelesti circulum maximum proiiciet, in quo locus heliocentricus terrae, locus heliocentricus corporis coelestis eiusdemque locus geocentricus siti erunt, et quidem ita ut distantia secundi a primo, tertii a secundo, tertii a primo, secundum eandem directionem numeratae, resp. sint $S$, $T$, $S+T$.

\SourceLeaf{084}IV. Vel ex notis variationibus differentialibus partium trianguli plani, vel aeque facile e formulis art. 62 sequentes aequationes differentiales derivantur:
\[
\mathrm{d}l=\frac{r'\cos(\lambda-l)}{\Delta'}\,\mathrm{d}\lambda+\frac{\sin(\lambda-l)}{\Delta'}\,\mathrm{d}r'
\]
\[
\mathrm{d}\Delta'=-r'\sin(\lambda-l)\,\mathrm{d}\lambda+\cos(\lambda-l)\,\mathrm{d}r'
\]
\[
\begin{aligned}
\mathrm{d}b={}&\frac{r'\cos b\sin b\sin(\lambda-l)}{\Delta'}\,\mathrm{d}\lambda
+\frac{r'\cos b^2}{\Delta'\cos\beta^2}\,\mathrm{d}\beta\\
&+\frac{\cos b^2}{\Delta'}(\operatorname{tang}\beta-\cos(\lambda-l)\operatorname{tang}b)\,\mathrm{d}r',
\end{aligned}
\]
\noindent ubi partes quae continent $\mathrm{d}r'$, $\mathrm{d}\Delta'$ per $206265^{\prime\prime}$ sunt multiplicandae, vel reliquae per $206265^{\prime\prime}$ dividendae, si mutationes angulorum in minutis secundis exprimuntur.

V. Problema inversum, scilicet determinatio loci heliocentrici e geocentrico, problemati supra evoluto prorsus analogum est, quamobrem superfluum foret, illi amplius inhaerere. Omnes enim formulae art. 62 etiam pro illo problemate valent, si modo, omnibus quantitatibus quae ad locum corporis coelestis heliocentricum spectant cum analogis iis quae ad geocentricum referuntur permutatis, pro $L$, $B$ resp. substituitur $L+180^0$, $-B$, sive quod idem est pro loco heliocentrico terrae geocentricus Solis accipitur.
